% !PRINTED 530
\weberpage{530}
%<<CONT
\[
\begin{aligned}
v_1 &= a_1^2 - 4(\alpha+\alpha_1)(\alpha_2+\alpha_3)\\
y_1-y_2 &= (\alpha+\alpha_1)(\alpha_2+\alpha_3) - 2(\alpha\alpha_1+\alpha_2\alpha_3)\\
a_2 &= (\alpha+\alpha_1)(\alpha_2+\alpha_3) + \alpha\alpha_1 + \alpha_2\alpha_3,
\end{aligned}
\]
also, wenn man $v_2$, $v_3$ ebenso bildet:
\weq{12}{
\begin{aligned}
3v_1 &= 3a_1^2 - 8a_2 - 4y_1 + 4y_2\\
3v_2 &= 3a_1^2 - 8a_2 - 4y_2 + 4y\\
3v_3 &= 3a_1^2 - 8a_2 - 4y + 4y_1.
\end{aligned}
}
Es kommt aber noch eine Relation hinzu, die sich daraus ergiebt, dass das Product
\[
(\alpha+\alpha_1-\alpha_2-\alpha_3)(\alpha-\alpha_1+\alpha_2-\alpha_3)(\alpha-\alpha_1-\alpha_2+\alpha_3)
\]
eine symmetrische Function der $\alpha$ ist, die man leicht gleich
\[
a_1^3 - 4a_1a_2 + 8a_3
\]
findet. Wir machen also die Vorzeichen der drei Wurzelgrössen durch die Relation
\weq{13}{\sqrt{v_1}\sqrt{v_2}\sqrt{v_3} = a_1^3 - 4a_1a_2 + 8a_3}
von einander abhängig und setzen:
\weq{14}{
\begin{aligned}
\sqrt{v_1} &= \alpha+\alpha_1-\alpha_2-\alpha_3\\
\sqrt{v_2} &= \alpha-\alpha_1+\alpha_2-\alpha_3\\
\sqrt{v_3} &= \alpha-\alpha_1-\alpha_2+\alpha_3\\
a_1 &= \alpha+\alpha_1+\alpha_2+\alpha_3.
\end{aligned}
}
Daraus ergiebt sich dann:
\weq{15}{
\begin{aligned}
4\alpha &= a_1+\sqrt{v_1}+\sqrt{v_2}+\sqrt{v_3}\\
4\alpha_1 &= a_1+\sqrt{v_1}-\sqrt{v_2}-\sqrt{v_3}\\
4\alpha_2 &= a_1-\sqrt{v_1}+\sqrt{v_2}-\sqrt{v_3}\\
4\alpha_3 &= a_1-\sqrt{v_1}-\sqrt{v_2}+\sqrt{v_3}.
\end{aligned}
}
Wenn man der $\sqrt{D}$ in (8) das entgegengesetzte Zeichen giebt, und $y$ durch $-y$ ersetzt, so muss man nach (10) $y_1$, $y_2$ mit $-y_2$, $-y_1$ vertauschen. Es bleibt also $v_1$ ungeändert und $v_2$ und $v_3$ werden mit einander vertauscht. Nimmt man für $y$ eine andere Wurzel der Gleichung (8), etwa $y_1$, so muss wegen (9) $y_1$, $y_2$ durch $y_2$, $y$ ersetzt werden und die Grössen $v_1$, $v_2$, $v_3$ erleiden eine cyklische Permutation, wodurch auch die Grössen $\alpha_1$, $\alpha_2$, $\alpha_3$ in (15) cyklisch vertauscht werden. Wenn man endlich die Vorzeichen von $\sqrt{v_1}$, $\sqrt{v_2}$, $\sqrt{v_3}$ anders wählt, aber so, dass immer (13) erfüllt bleibt, so treten unter den Grössen (15) gleichfalls gewisse Permutationen ein.
